\section{Discussion}

\subsection{AdsMind's Contributions and Scope}

Results point to a nuanced characterization of closed-loop LLM agents.
Prior work emphasizes \textit{capability} (materials discovery, autonomous simulation, catalyst screening); for finding low-energy configurations, capability is not where closed-loop reasoning wins.

\textbf{Energy-depth comparison.}
Random sampling (20 relaxations) matches AdsMind on 10/20 CMU cases, beats it on six.
Heuristic enumeration exposes additional low-energy basins but requires connectivity filtering.
Adsorb-Agent's multi-configuration search reaches lower energies in 14/16 (GPT-5.4) and 12/16 (GPT-4o) comparable cases.

\textbf{Robustness across backends.}
Results hold on CMU20 and OCD-GMAE62 (two-tier protocol) across four backends after removing external service failures.

Closed-loop reasoning contributes along a different axis: \textit{reliable, compute-efficient, and interpretable} simultaneously.
Each attribute has precedents alone; AdsMind's novelty is their \textit{co-occurrence}:
\textbf{Reliability:} 80/80 CMU20 Full; OCD-GMAE62 Tier~1: 40/40 Full vs.\ 36/40 1-Shot; 320/320 iterative vs.\ 73/80 1-Shot.
\textbf{Efficiency:} 5--25$\times$ less MACE compute than brute-force.
\textbf{Interpretability:} physically grounded artifacts (slip events, forbid constraints, site-mismatch fingerprints) for inspection~\cite{zhang2026matclaw,huang2026towards}.
This combination is prerequisite for autonomous chemistry; recent work emphasizes similar requirements~\cite{zhang2026towards,yang2026quasar}.



\subsection{Claims and evaluation}

Four empirical claims established through ablation attempts (400 CMU20 + OCD-GMAE62 Tier~1 1240 + Tier~2 720 repeated runs), Tier~1 generalization study, and matched baselines:

\textbf{C$_1$ (reliability):} Full achieves 100\% on CMU20 (80/80) and OCD-GMAE62 Tier~1 (245/248, 98.8\%).
CMU20: 320/320 iterative vs.\ 73/80 1-Shot.
OCD-GMAE62 Tier~2 (12 cases, N=3): stability audit shows 131/240 comparisons within 0.01~eV across runs.
Adsorb-Agent: 73--80\% success; single-configuration stress test fails on 6/20.
All shortfalls are natural outcomes, not service failures.

\textbf{C$_2$ (backend convergence):} Iterative architecture reduces energy variance.
CMU20: Full 0.153~eV vs.\ 1-Shot 0.359~eV (2.3$\times$ reduction).
OCD-GMAE62 Tier~1: Full 0.183~eV.
OCD-GMAE62 Tier~2 (stability): 27/47 Full runs within 0.01~eV across 3 repeats; mean range 0.166~eV.

\textbf{C$_3$ (compute efficiency):} Single Full run uses $\sim$4 relaxations.
Comparators: 20 (random), 25--98 (heuristic), $\sim$21 (Adsorb-Agent).
CMU20: AdsMind matches/beats random on 14/20 (0.01~eV tolerance).
Heuristic: 11 ties, 1 AdsMind-lower, 3 heuristic-lower (5 unresolved).

\textbf{C$_4$ (interpretability):} Closed-loop reasoning produces structured artifacts:
slip events, forbid constraints, site-mismatch fingerprints.
One-shot slip rates correlate with surface complexity.
Chemical Slip serves as standalone LLM evaluation signal.


\subsection{Chemical Slip as evaluation signal}

Chemical Slip (mismatch between LLM-planned and realized adsorption site) is a feedback signal for closed-loop agents.
Slip statistics from CMU20 and OCD-GMAE62 reveal secondary value: slip frequency (stratified by surface type/adsorbate) is a physically grounded proxy for LLM chemical reasoning quality, independent of downstream metrics~\cite{alfred2026data,zhang2026towards}.
One-shot slip rates: 20\% (monometallic) to 77\% (intermetallic), indicating LLM priors are less reliable on complex alloys.
Unlike text-based benchmarks (ChemBench~\cite{mirza2024chembench}), Chemical Slip evaluates \textit{applied} reasoning (whether structural prediction survives physical simulation) and complements existing LLM evaluation frameworks~\cite{li2026agentic,huang2026towards}.
Chemical Slip is a general-purpose evaluation signal for LLM scientific agents.

\subsection{Why brute-force sampling wins on depth, and when it matters}

Energy-depth advantage of random/heuristic methods requires mechanistic interpretation.
On simple adsorbates (few binding sites), AdsMind Full matches non-LLM baselines within 0.01~eV.
On complex adsorbates (CH$_2$CH$_2$OH, etc.), PES has many local minima; iterative refinement from single trajectory becomes trapped~\cite{lan2023adsorbml,zitnick2020introduction}.

This delineates when closed-loop search is the right tool.
For high-throughput screening with abundant parallel compute, brute-force can be more wall-clock efficient (failures tolerable, though $5$--$25\times$ more relaxations).
For autonomous exploration of single candidate where failure is costly, closed-loop adds value.
Two-stage pipeline (brute-force + closed-loop refinement) would combine throughput with reliability.

\subsection{MLFF as surrogate: scope and validation}

All experiments use MACE-MP-0 small as physical backend~\cite{batatia2023foundation,deng2023chgnet}.
The agent's value is orthogonal to MLFF accuracy: it optimizes search strategy, not energy function.
Adsorption energy is operationally defined within fixed MACE-MP-0 small protocol; within-protocol rankings are self-consistent.
We do not claim MACE-MP-0 small rankings match DFT rankings; MACE-M-P-0 large serves as sensitivity boundary, not ground-truth validation.

We repeated GPT-5.4 Full on all 20 CMU cases with MACE-MP-0 large (GPU, float64, dispersion enabled).
Mean absolute shift vs.\ MACE-MP-0 small: 1.255~eV (0.834~eV excluding case~20); 15/20 more stable, 5/20 less stable.
Case~20: retained non-dissociated structure at $-2.402$~eV; later dissociated analyses reach lower unretained energies.
On OCD-GMAE62 Tier~1 sensitivity subset: MACE-M-P-0 large succeeds on 9/10 cases; case~043 fails naturally.
This is a force-field sensitivity check, not AdsMind--Adsorb-Agent head-to-head.
DFT validation remains necessary for final energetics beyond fixed surrogate protocol.

Known MLFF failure modes---dispersion-dominated binding, strongly correlated $f$-electron systems, strong charge-transfer interfaces---define the current scope boundary of the framework.
Extending to these regimes will require either higher-fidelity MLFFs or hybrid MLFF--DFT workflows, but does not require changes to the agent architecture.

\subsection{Limitations and negative results}

\textbf{Algorithmic.} On complex adsorbates (16, 17, 18, 19), AdsMind Full produces higher energies than broader-sampling baselines.
Five-seed control reduces but does not remove gap (0.90~eV on case~16, 1.20~eV on case~17).
Mechanism: iterative trajectory stays within initial basin; this is specification, not defect.

FORBID over-pruning: removing constraint \textit{improves} energy on several cases (GPT-5.4: BH $p = 0.084$).
Indicates explicit negative constraints over-prune; adaptive variants merit future investigation.

\textbf{Coverage.} Exhaustive controls show dominant factor is \textit{number of initial configurations sampled}.
Random (20 relaxations) finds lower energies on 6/20; heuristic (25--98 sites) wins 3/20; Adsorb-Agent ($\sim$21) wins 14/16 (GPT-5.4), 12/16 (GPT-4o).
Broader sampling matches/exceeds closed-loop refinement at 5--25$\times$ more compute.
Question: what is contribution of closed-loop reasoning when compute/physics/data are fixed?

\textbf{Statistical.} Friedman tests reach significance ($p < 0.01$); many Wilcoxon tests do not survive BH correction.
We report failure-aware rates and success-conditioned deltas; interpret via effect sizes and bootstrap intervals.

\textbf{Scope.} Current implementation: flat/stepped crystalline surfaces, single adsorbate.
Extensions (defects, nanoparticles, coadsorption, solvation) are architecturaly straightforward but require validation~\cite{yang2026quasar,zhang2026towards}.

\textbf{LLM reproducibility.} Hosted services are black-box, subject to unannounced updates.
We mitigate by recording model identifiers, frozen settings, releasing result files and scripts.
Multi-backend design ensures conclusions do not depend on single backend~\cite{alfred2026data,huang2026towards}.


\section{Conclusion}

We introduced AdsMind, a closed-loop multi-agent framework for autonomous adsorption configuration search that integrates LLM-driven planning with iterative MLFF relaxation.
By framing adsorption configuration search as a sequential decision problem, AdsMind overcomes the fundamental limitation of open-loop LLM agents: the inability to correct reasoning errors based on physical feedback.
Three mechanisms---Chemical Slip detection, FORBID constraint memory, and autonomous termination---enable the agent to detect, diagnose, and recover from incorrect placement hypotheses.
Through 880 ablation attempts across four LLM backends on two benchmarks (CMU20 and OCD-GMAE62), plus head-to-head comparisons with existing methods, we established that AdsMind's primary value lies in four properties that co-occur: \textbf{Reliability}: 100\% success on CMU20 (80/80) and 98.8\% on OCD-GMAE62 (245/248), versus 91.3\% and 89.5\% for 1-Shot; \textbf{Backend convergence}: 2.3$\times$ reduction in cross-backend energy variance (0.153~eV vs.\ 0.359~eV), demonstrating that iterative feedback suppresses LLM-specific placement bias; \textbf{Compute efficiency}: $\sim$4 relaxations per case versus 20--98 for brute-force baselines (5--25$\times$ reduction); \textbf{Interpretability}: structured artifacts (slip events, forbid constraints, site-mismatch fingerprints) provide physically grounded inspection trails, while Chemical Slip serves as a standalone LLM evaluation signal.
A central finding is that closed-loop reasoning does not replace brute-force sampling for energy-depth search; rather, it provides a different axis of improvement---reliability, efficiency, and interpretability simultaneously---that brute-force methods cannot match.
On complex adsorbates with many local minima, broader initial sampling reaches lower energy basins when it succeeds; however, it fails silently on 8.7--10.5\% of cases and lacks chemical-safety filters.
AdsMind's architectural contribution is decoupling \textit{reliable convergence} from \textit{exhaustive sampling}, enabling autonomous workflows where failure is costly or parallel compute is limited.
Future work should couple AdsMind's closed-loop recovery with broader enumeration strategies, inheriting the depth of brute-force sampling and the reliability of iterative refinement.
Extending to co-adsorption, solvation, and defected surfaces requires minimal architectural changes but significant validation against DFT reference data.
More broadly, the Chemical Slip diagnostic---measuring mismatch between LLM-planned and physically realized binding sites---provides a general-purpose evaluation signal for LLM-based scientific agents beyond surface science.
As LLMs increasingly participate in scientific discovery, frameworks like AdsMind that ground language-model reasoning in physical feedback will be essential for robust, interpretable autonomy.
